\section{Experimental Setup}
\label{sec:setup}

\paragraph{플랫폼.}
실물 기준 플랫폼은 ROBOTIS AI Worker FFW-SG2다(이동 베이스 고정). 7자유도 팔 두 개와 평행 그리퍼, 2자유도 머리, 리프트를 쓰고 팔마다 100\,Hz 위치 명령으로 제어한다. 카메라는 기본 장착분만 쓴다. 머리는 ZED Mini 왼쪽 정류 영상 672$\times$376(수평 85$^\circ$), 손목은 RealSense D405 424$\times$240 두 대다. 손목 힘/토크 센서는 없다.

\paragraph{시뮬레이션과 과제.}
Isaac Sim 5.1과 Isaac Lab 2.3에서 공개 챌린지 환경의 로봇·카메라 설정을 그대로 쓴다. 개발 과제는 오른팔 탁상 집어 놓기(``Put the red mug on the blue tray.'')이며, 성공은 머그가 트레이 위에 바로 서 있고 쥐고 있지 않은 상태가 1\,s 이어지는 것이다. 학습 데이터 생성기는 세 과제(머그를 트레이에, 머그를 표적에, 병을 트레이에)를 지원한다. 물리는 CPU PhysX다. 같은 시드가 직전 실행 이력에 따라 다른 궤적으로 재생되는 문제를 찾아, 매 에피소드 시작에서 물리 장면을 다시 만들도록 고쳤다(서로 다른 이력 뒤의 실행이 매 틱 비트 단위로 같음; 부록 \cref{sec:supp_repro}).

\paragraph{장면과 분할.}
\textbf{standard}는 기본 장면이다. \textbf{random}은 방해물·탁자·바닥 재질·조명·배경의 5축을 시험 전용 풀에서 뽑은 장면이고, \textbf{DR}은 같은 5축을 겹치지 않는 학습용 풀에서 뽑는다. 시드는 용도별로 나눈다. 개발 DEV(0--29), 보정 CAL, 시험 TEST, 학습 R2\_TRAIN(10000번대), 그리고 학습 금지 OOD 분할(70000--70999)이다. 보호 분할은 해당 실험의 사전 등록 뒤에만 열린다. 학습 데이터 R2\_TRAIN은 6,000편을 생성해 유효 \tentative{5,126}편이다.

\paragraph{평가 하네스와 시계.}
1차 평가는 오픈소스 하네스 Inspect Robots~\cite{inspectrobots2026}에서 한다. 하네스 루프는 동기식이라 시계만 바꿔 넣는다. \textbf{지연 충실 트랙}(주 표)에서는 요청의 응답을 실측 지연 뒤에 받고, 동기식 기준선은 응답을 계산하는 동안 직전 명령을 유지한 것으로 처리한다. \textbf{동기 공정 트랙}은 생각 시간을 공짜로 두며 같은 표에 함께 적는다. 비정지 실행의 이점은 지연 충실 트랙에서만 주장한다.

\paragraph{비교 대상.}
(1) 상위 계획기 단독(동기, Astra에 맞춘 인터페이스), (2) VLA 단독, (3) 학습 VLA($\pi_{0.5}$ 미세조정), (4) 같은 상태 위의 코드 규칙과 스크립트 스킬, (5) 코드 정책~\cite{fu2026capx}, 그리고 우리 결합 계층의 변형(상위 = Astra, 학습한 8B·35B 학생). 모든 비교는 같은 layout과 seed, 같은 정보(과제 문장만; 오라클 물체 자세는 주지 않음), 같은 상위 호출 예산, 고정한 모델 판본으로 한다.

\paragraph{통계와 사전 등록.}
비율 차이의 95\% 구간은 같은 시드 짝에 대한 군집 부트스트랩 10,000회이고, 여러 조건이 걸리면 Holm 보정을 한다. 모든 실험은 실행 전에 ``이 결과로 바꾸는 설계 결정'', ``이 표본으로 그 결정을 가를 수 있는가'', 판정 문턱을 적어 커밋한다. 무료 로컬 모델로 먼저 돌려 배선 결함을 잡고, 관문에서 결함이 보이면 멈추고 다시 설계한다. 판정 스크립트는 결과 전에 고정하고, 독립 검증자가 등록 문서와 코드를 한 줄씩 대조하는 적합성 감사를 두 번 연속 통과해야 본 실행에 들어간다(부록 \cref{sec:supp_repro}).

\paragraph{계산 자원.}
시뮬 렌더, 학습, 추론 검증은 H200 파드에서 한다. 실물에서는 로봇 옆 GPU 서버(RTX PRO 6000)를 가정하며, 지연 수치는 실물 서버에서 다시 잰다. 상위 모델 Astra는 API로 부른다.
